\chapter{Breast Tenderness (Mastalgia)}

\section*{Definition}
Pain, discomfort, or tenderness in one or both breasts, ranging from mild soreness to severe pain that interferes with daily activities, often cyclical in relation to hormonal fluctuations.

\section*{What Happens in Your Body}
Cyclical mastalgia correlates with the luteal phase of the menstrual cycle when progesterone and oestrogen peak, causing breast tissue oedema, ductal proliferation, and stromal expansion. Hormone-sensitive breast tissue contains oestrogen and progesterone receptors; hormonal fluctuations during perimenopause alter the normal cyclical pattern. Non-cyclical mastalgia arises from extramammary structures (chest wall, costochondral junctions) or focal breast pathology (cysts, fibroadenomas, duct ectasia).

\section*{Causes}
\begin{itemize}
 \item Cyclical hormonal fluctuations (perimenopause, menstrual cycle)
 \item Hormone replacement therapy (HRT) or oral contraceptives
 \item Breast cysts (fibrocystic breast changes)
 \item Fibroadenomas
 \item Mastitis (infectious or non-infectious inflammation)
 \item Duct ectasia
 \item Costochondritis (chest wall pain referred to breast)
 \item Trauma or prior breast surgery
 \item Large or heavy breasts (mechanical strain)
 \item Breast cancer (rare, usually painless but possible)
\end{itemize}

\section*{When to See a Doctor}

\begin{itemize}
  \item Symptoms are severe or getting worse over time
 \item Breast pain with a new palpable lump, skin changes (dimpling, redness, peau d'orange), nipple discharge, or persistent focal pain lasting more than one menstrual cycle
  \item You have other concerning symptoms such as fever, weight loss, or bleeding
\end{itemize}

\section*{Related Symptoms}
\begin{itemize}
 \item Breast lump
 \item Nipple discharge
 \item Breast swelling
 \item Premenstrual syndrome
 \item Painful sex
\end{itemize}

\textbf{Important:} Your doctor will check for serious underlying conditions when evaluating persistent breast tenderness.

\section*{Self-Care / Home Management}
\begin{itemize}
 \item Wear a well-fitted, supportive bra with wide straps and avoid underwire bras during symptomatic periods.
 \item Apply cold packs to the breasts for 15–20 minutes to reduce discomfort during tender days.
 \item Reduce dietary caffeine and sodium intake in the week before the menstrual cycle, which may help cyclical mastalgia.
 \item Take over-the-counter NSAIDs (ibuprofen, naproxen) for pain relief, starting a few days before the expected onset of symptoms if cyclical.
\end{itemize}

\section*{How Your Doctor Will Evaluate You}
\begin{itemize}
 \item Your doctor will differentiate cyclical from non-cyclical mastalgia by history: timing relative to menses, location, and relationship to hormonal therapy.
 \item Clinical breast examination is performed to identify any focal tenderness, masses, or skin changes that require further imaging.
 \item Bilateral diagnostic mammogram or breast ultrasound is arranged for unilateral, focal, or persistent breast pain to exclude underlying pathology.
 \item Referral to a breast specialist is indicated for focal non-cyclical pain with a suspicious imaging finding or if symptoms significantly impair quality of life.
\end{itemize}

\section*{Treatment Options}
\begin{itemize}
 \item Cyclical mastalgia is managed with lifestyle measures (supportive bra, reduced caffeine and sodium) and NSAIDs, with danazol or tamoxifen reserved for severe cases.
 \item Adjustment or switching of hormonal therapy (HRT, oral contraceptives) may reduce breast tenderness if the current regimen is a contributing factor.
 \item Evening primrose oil or vitamin E supplementation may provide modest benefit for cyclical mastalgia.
 \item Treatment of the underlying cause includes antibiotics for mastitis, aspiration of symptomatic cysts, and surgical excision for fibroadenomas or other focal lesions.
 \item Medications, lifestyle modifications, and physical therapies may be prescribed as appropriate.
 \item Follow-up ensures treatment effectiveness and allows timely adjustments to the care plan.
\end{itemize}

\section*{References}
\begin{thebibliography}{99}
\bibitem{ref1} Salzman B, Collins E, et al. "Common breast problems." *Am Fam Physician*. 2019;99(7):446-453.
\bibitem{ref2} Smith RL, Pruthi S, et al. "Evaluation and management of breast pain." *Mayo Clin Proc*. 2004;79(3):353-358.
\bibitem{ref3} Rosolowich V, Satin N, et al. "Mastalgia: a review." *J Obstet Gynaecol Can*. 2006;28(1):49-57.
\bibitem{ref4} Harris JR, Lippman ME, et al. *Diseases of the Breast*, 5th ed. Wolters Kluwer, 2014.
\bibitem{ref5} UpToDate. "Breast pain: evaluation and management." Wolters Kluwer, 2023.
\end{thebibliography}
